% !TEX root = ../../main/main.tex
\hypearticle
{IA na segurança: o impacto no cidadão comum e o futuro da sociedade}
{Luiz Claudio Franco}
{

Todos nós estamos cientes do recente crescimento da integração de inteligência artificial em diversas áreas, a fim de otimizar processos e automatizar tarefas que até então eram realizadas mecanicamente por humanos. Até aí tudo certo, nada disso é novidade.\par

Desde os primórdios da nossa história, a tecnologia desempenha a função de criar soluções cada vez mais eficazes para o cotidiano do homem. Por isso, imagino que o leitor concorde que não há nada de errado no uso da IA para aperfeiçoar nossas tarefas, certo? Bom, a resposta para esse dilema é \textbf{tudo menos binária}.\par

Recentemente, a polícia dos Estados Unidos tem adotado sistemas de reconhecimento facial baseados em IA para identificar fugitivos e suspeitos de crimes. No entanto, essa tecnologia vem produzindo inúmeros \textbf{falsos positivos}, resultando na prisão de pessoas inocentes. Um incidente que chama a atenção é o de uma senhora de Oklahoma que foi presa injustamente por seis meses por conta de um pedido de prisão vindo de Maryland. Qual a parte mais absurda disso? A mulher nunca sequer esteve no Estado de onde o pedido saiu.\par

Ao analisar esse caso, dois fatores chamam a atenção: a \textbf{incapacidade da polícia} de apurar os resultados dados pelo sistema de reconhecimento, e a \textbf{falta de vontade do tribunal} em analisar as evidências apresentadas pelo réu, que provavam claramente que ela não poderia estar naquele Estado na data do suposto crime cometido.\par

Outro caso alarmante ocorreu no Reino Unido, onde um detetive da polícia de Derbyshire utilizava inteligência artificial para \textbf{forjar evidências} em diversos processos.\par

As consequências desse tipo de acontecimento são \textbf{devastadoras}. No caso de Oklahoma, a mulher foi privada de liberdade por seis meses, além de ter perdido seu emprego e de ter seu histórico manchado. Imagine se esse tipo de situação se tornasse cada vez mais comum, em que qualquer pessoa inocente poderia ser encarcerada por um simples erro de IA ou pela malícia de um agente público manipulando o sistema. Este cenário descreve um \textbf{futuro distópico e sombrio} para a humanidade.\par

Cada um dos ocorridos apresentados revela os dois aspectos mais perigosos da IA na atualidade: a \textbf{confiança cega} nos resultados trazidos por ela, e a capacidade de gerar materiais artificiais que podem ser interpretados como legítimos \textbf{à primeira vista}.\par

Logo, chegamos à raiz do problema: \textbf{o ser humano}. O maior obstáculo para usufruir da tecnologia de maneira responsável nunca foi e nunca será a ferramenta em si, mas sim \textbf{quem a opera}. A solução não é colocar todas as IAs na mesma caixinha e dizer que são prejudiciais; isso não seria justo nem benéfico para a sociedade. Até mesmo os impactos ambientais que esses datacenters vêm trazendo são apenas o produto do uso exagerado da inteligência artificial pelo homem, muitas vezes para gerar conteúdos banais e irrelevantes (como memes e afins).\par

Como prosseguir, então? É necessário \textbf{educar a sociedade} para o uso correto e responsável da tecnologia, utilizando-a como um auxílio e não como a dona da verdade. Infelizmente, o passado nos mostra que o homem tende a abusar das inovações, seja para usá-las contra o outro, seja como muleta. Portanto, a evolução tecnológica não é apenas um caminho unidirecional, mas também um processo constante do ser humano para promover um uso cada vez mais \textbf{responsável e ético}.\par

\highlight{
    \textcolor{HypePurple}{Seja responsável com a tecnologia!}\\
}

\originalsource{https://www.aclu.org/news/privacy-technology/more-than-a-dozen-wrongful-arrests-due-to-police-reliance-on-facial-recognition-technology}

\textcolor{HypePurple2}{Outras fontes:}\\
\href{https://www.bbc.com/news/articles/c7vgq5mgz2mo}{BBC — caso de uso de IA em investigação}


}